\section{Independent checks of comparison, identification and class limits}
\subsection{Paired inference conditional on the number of discordances}
For a frozen pair of policies let $A,B\in\{0,1\}$ denote losses on the same mission. Let $(N_{10},N_{01})$ be the counts of the discordant pairs in $n>0$ independent copies of the joint experiment. Write $\theta=\Pr(A\ne B)$ and, when $\theta>0$, $\varphi=\Pr(A=1,B=0\mid A\ne B)$. Then the risk difference is $\Delta=\theta(2\varphi-1)$. Conditional on $M=N_{10}+N_{01}=m$, $N_{10}\sim{\rm Bin}(m,\varphi)$, whereas $M\sim{\rm Bin}(n,\theta)$. To see this, aggregate the two concordant outcomes into a third multinomial cell. Each multinomial probability factors into the binomial probability of $m$ discordances and a binomial probability of their orientation. This factorization allows arbitrary dependence within the pair.

Use any exact binomial interval at total error probability $\gamma/2$ for $\theta$ and, when $m>0$, a second such interval for $\varphi$ given $m$. For $m=0$ use $[0,1]$ for $\varphi$. Conditional orientation noncoverage is at most $\gamma/2$ for every $m$, so its unconditional noncoverage is at most $\gamma/2$. The union inequality therefore covers both parameters at level $1-\gamma$. Since the map $(\theta,\varphi)\mapsto\Delta$ is affine in each argument, its extrema on the resulting rectangle occur at vertices. At $\theta=0$ the truth is $\Delta=0$ regardless of how orientation is defined. This construction and the main benchmark's separate binomial bounds on the two discordant cell probabilities are two valid procedures. They need not give identical endpoints. Looking at both and taking the narrower interval requires a shared error allocation.

Zero observed discordances do not imply identical future behavior. For example, let $B=0$ always and $A=1$ with probability $p>0$. Then pairs match on every one of $n$ independent test missions with probability $(1-p)^n$. Proving that two deterministic algorithms have identical outputs for all admissible inputs is different evidence. Conversely, for pathwise identical loss code the contrast is exactly zero even when the future event rate itself is uncertain.

\subsection{Why an uncertain competitor enters a regret bound}
Take a required false-alarm probability at most $0.04$. For policy $A$, suppose a simultaneous enclosure of false alarm and hazard risk is $[.01,.03]\times[.20,.30]$. For $B$ it is $[.01,.06]\times[.01,.05]$. The validation evidence certifies $A$ as feasible. Two compatible worlds have $p_{{\rm FA},B}=.06$ or $.01$ and the same true hazard risks $.30$ for $A$ and $.01$ for $B$. In the former, $A$ is the only feasible option. In the latter, selecting $A$ has excess risk $.29$. No universal guarantee smaller than $.29$ follows from these rectangles alone. A paired contrast or additional information on feasibility can lower this value. Checking only certified competitors would incorrectly return zero here. The generic engine is tested independently against all $2^6$ endpoint worlds for three policies, at five budgets including equality at the boundary. For each world it checks feasibility containment, optimum enclosure and the reported regret inequality against the directly computed best feasible candidate. These are computational cross-checks of the proof and do not replace its quantification over all interior values.

\subsection{Partial identification with common latency}
Consider a population in which every mission requires an intervention. Record the paired actionable margins, under a common untreated state history before any alarm-induced physical change. For an ordered latency support $r_1<\cdots<r_K$, encode a margin by the number of listed durations that it covers, with zero for an unusable alarm. For $m$ policies, let $a$ be the $m$-component vector of these counts, $p_a$ its mission probability, and $q_j$ the probability of latency $r_j$. An exogenous common actuator latency means that each realization of the latency is applied to all counterfactual policies on that mission. It need not be statistically independent of the margins. Write $z_{aj}$ for their unknown joint probabilities. Given exact marginals, every possible common law induces a nonnegative matrix with row sums $p_a$ and column sums $q_j$. Conversely every such matrix defines a common admissible law: draw $(a,j)$ with mass $z_{aj}$ and select observations from the conditional mission distribution for type $a$. Under the model specifying only type and latency marginals, this attains the corresponding risk or contrast.

For a pair $(i,h)$ the cell contrast is $\mathbf1\{a_i<j\}-\mathbf1\{a_h<j\}$. Minimizing and maximizing its linear expectation over the transport polytope therefore gives a sharp identification interval. There is no requirement that two coordinates attain their individual maxima on the same law. For instance, two policies with identical type on every mission always have contrast zero. If both the short/long margin type and short/long latency have probability $1/2$ each, identical marginals may allow individual failure from zero to $1/2$ under alternative couplings. Subtracting separate probability intervals can then give $[-1/2,1/2]$ even though the true contrast equals zero under every common coupling.

With independently observed complete missions, form binomial intervals for each type probability, including types with zero counts, with simultaneous coverage via a union bound. If actuator shares are also sampled, include their interval failures in the same budget. Optimize over unit-mass matrices with row and column probabilities in those intervals. On simultaneous coverage the actual joint probability matrix is feasible, so the extrema enclose each admissible contrast at once. This construction is sharp for fixed exact marginals; enlarging to a confidence rectangle is conservative. Mathematically unattainable types may be removed only with a justification independent of their absence from the sample. This model does not address feedback that makes latency policy-specific. In that setting a common scalar $r_j$ imposes an additional, empirically consequential coupling assumption.

\subsection{Independent verification of the finite causal frontier}
The three-epoch sensor example admits a third numerical method in addition to rational policy enumeration and the backward support calculation. Let $\mathcal H$ be its 14 nonempty binary histories and $\mathcal L$ its eight length-three leaves. Introduce $x_h$ as the probability with respect to the monitor's independent random bits that it first stops at history $h$, conditional on that exact exogenous sensor history. Every nonanticipating randomized rule satisfies $x_h\ge0$ and
\[
 \sum_{h:\,h\preceq \ell} x_h\le1,\qquad \ell\in\mathcal L.
\]
Here $h\preceq\ell$ means a prefix. The rule's randomization is independent of future symbols, hence the same prefix probabilities apply to every continuation. Conversely choose any such vector. At a reached node $h$, given no previous stop, halt with probability
\[
 \frac{x_h}{1-\sum_{v:\,v\prec h} x_v}.
\]
When the denominator is zero the node is unreachable under continued operation and $x_h=0$; its continuation can be set arbitrarily. Induction gives the prescribed prefix probabilities. This establishes exactness of the linear constraints, including all randomized policies. It uses exogenous observations, a single terminal action, no side information and the same randomization law in the nominal and changed hypotheses.

For a history of length $t$ with $k$ ones define $b_p(h)=p^k(1-p)^{t-k}$. For the stated mixture with early/late deadline probabilities $.4$ and $.6$, the coefficient of timely warning is
\[
 c_h=.4\,\mathbf1\{t=1\} b_{.8}(h)+.6\,b_{.6}(h).
\]
The feasible false-alarm constraint is $\sum_h b_{.2}(h)x_h\le\alpha$. Maximizing $\sum_h c_hx_h$ over the prefix constraints is thus exactly the constrained causal value. The accompanying test solves this real-valued LP independently of the rational stopping-tree code and checks each published budget. At $\alpha=.10$, the optimal timely probability is $59/150$, corresponding to hazard probability $1-.9(59/150)=323/500$ under this example's independent intervention efficacy. The full-history terminal test bound is $.73$ and does not give this causal optimum. Rational enumeration and exact support-function equalities establish the discrete result without depending on a floating-point LP tolerance. No precision or optimality conclusion carries over to a state-dependent intervention response with a different information structure.
